% ============================================================================
%  Chapter 6: Conclusion
% ============================================================================
\section{Conclusion}
\label{sec:conclusion}

We presented a decentralized guarantee and arbitration layer for agent-to-agent commerce built on
the Schelling-point reputation community: registered agent identities, a guaranteed escrow,
staked-voting adjudication, and reputation-driven guarantee pricing, closed into one on-chain loop.
We proved that under a two-thirds supermajority rule honest voting is incentive compatible, that a
coordinated malicious minority of at most $35\%$ of voters is strategically harmless, that Sybil
attacks are unprofitable because adversarial voting power scales linearly in expenditure, and that
the mechanism conserves funds exactly. We instantiated the protocol in Solidity with 38 passing
tests and an end-to-end demo, and validated the theory by Monte-Carlo simulation built on
\textsf{mesa}. To our knowledge this is the first protocol to tie adjudication, guarantees, and
reputation into a single verifiable cycle, addressing the accountability vacuum of anonymous
machine-to-machine commerce (Q2) with an economically self-sustaining governance mechanism (Q3).

\paragraph{Limitations.}
The current MVP makes several deliberate simplifications, each flagged on-chain and in the theory:
(i) \emph{free abstention} --- abstainers are refunded rather than penalized, so the one-shot
monetary payoff of honest voting is only marginally better than abstaining; the reputational
tie-breaker of the closed loop is therefore load-bearing for participation (Section~\ref{sec:pricing});
(ii) \emph{integer dust} --- the majority reward's division remainder is stranded in the contract;
(iii) \emph{engineering threshold} --- $\texttt{MAJORITY\_BPS} = 6500$ approximates two-thirds with
integer-safe arithmetic; (iv) \emph{demonstration chain} --- the deployment targets a local anvil
chain and, as a next step, a public testnet (Section~\ref{sec:deploy}); and (v) \emph{no random jury
selection} --- anyone can vote on any case, which the full design refines with ZK-augmented random
selection.

\paragraph{Future work.}
Short-term directions are: batch settlement to eliminate dust; the full identity barrier that binds
every vote to a registered, fee-paid agent ID; ZK-augmented random jury selection with privacy
adjustability; and a larger simulation study of threshold tuning, collusion dynamics, and
reputation-strategy equilibria. Longer term, the mechanism is designed to operate on a public
Ethereum L2 (Base) and to integrate with the ERC-8004 trustless-agent ecosystem; a production
deployment will follow the testnet checklist of Section~\ref{sec:deploy}. Compliance-wise, the
MVP uses testnet tokens with no real value; any tokenization would require an overseas-compliant
issuance structure, while liability remains anchored to real principals via the registry.

\paragraph{Concluding remark.}
The core design bet of this paper is that in an economy of autonomous agents, trust should be
\emph{earned, verified, and priced} --- not assumed. By making the community that arbitrates a
dispute the same community whose reputational and financial stakes depend on the verdict, the
Schelling-point reputation community turns truthful adjudication into the rational choice, and
turns reputation into a storable, revenue-bearing asset that prices the very guarantees that make
future agent-to-agent commerce safe.
